\subsubsection{The formal two-dimensional holomorphic OPE}
\label{sssec:formal-two-dimensional-holomorphic-ope}

The holomorphic OPE of the local model is a diagonal expansion on
\(\C^2\times\C^2\), not a Laurent expansion on a curve.  Let
\[
  \Delta\subset \C^2_z\times\C^2_w,\qquad
  \zeta_i=z_i-w_i ,
\]
and let \(j\colon(\C^2\times\C^2)\setminus\Delta\hookrightarrow
\C^2\times\C^2\) be the open complement.  The singular coefficient
module of a two-dimensional holomorphic OPE is the local-cohomology
module
\[
  \mathcal P_\Delta
  =
  H^2_\Delta\bigl(\C[[z_1,z_2,w_1,w_2]]\bigr)\,d\zeta_1d\zeta_2
  =
  \bigoplus_{a,b\ge0}
  \C[[w_1,w_2]]\,\rho^\Delta_{a,b},
\]
with
\[
  \rho^\Delta_{a,b}
  =
  \zeta_1^{-a-1}\zeta_2^{-b-1}\,d\zeta_1d\zeta_2,\qquad
  \operatorname{Res}_\Delta\bigl(\rho^\Delta_{a,b}\,
  \zeta_1^m\zeta_2^n\bigr)=\delta_{a,m}\delta_{b,n}.
\]
The two-dimensional Cauchy kernel is the generator
\[
  K^{(2)}_\Delta(z,w)=\rho^\Delta_{0,0}
  =
  \frac{d\zeta_1d\zeta_2}{\zeta_1\zeta_2}.
\]
In the region \(|z_i|>|w_i|\) it has the formal expansion
\[
  K^{(2)}_\Delta(z,w)
  =
  \sum_{a,b\ge0}
  z_1^{-a-1}z_2^{-b-1}w_1^aw_2^b\,dz_1dz_2 .
\]
This is the two-variable mode expansion used below.  A chiral algebra
on a curve is obtained only after replacing \(\mathcal P_\Delta\) by
the one-variable principal parts of a chosen line \(C\hookrightarrow\C^2\).

For \(x\in\mathfrak g=\mathfrak h\ltimes
 \mathfrak h^\vee_{\mathrm{cont}}[1]\), write \(\mathsf{J}_x(z)\) for
the linear holomorphic current determined by \(x\) in the CE
factorization algebra
\[
  B\longmapsto
  C^\bullet_{\mathrm{CE},\mathrm{cont}}
  \bigl(\Omega_c^{0,\bullet}(B)\widehat\otimes\mathfrak g\bigr).
\]
For \(f\in\mathfrak h\) and
\(\rho\in\mathfrak h^\vee_{\mathrm{cont}}\) we also write
\[
  \mathsf{B}_f=\mathsf{J}_f,\qquad
  \mathsf{\Theta}_\rho=\mathsf{J}_{\rho[1]} .
\]

\begin{thm}[Formal two-dimensional holomorphic OPE]
\label{thm:formal-two-dimensional-holomorphic-ope}
The singular part of the holomorphic OPE of the Hamiltonian BF
factorization algebra on \(\widehat{\C^2}_0\) is
\[
  \mathsf{J}_x(z)\mathsf{J}_y(w)
  \sim
  K^{(2)}_\Delta(z,w)\,\mathsf{J}_{[x,y]_{\mathfrak g}}(w),
  \qquad x,y\in\mathfrak g .
\]
Equivalently, for Hamiltonians \(f,g\in\mathfrak h\) and principal
parts \(\rho,\sigma\in\mathfrak h^\vee_{\mathrm{cont}}\),
\[
\begin{aligned}
  \mathsf{B}_f(z)\mathsf{B}_g(w)
  &\sim
  K^{(2)}_\Delta(z,w)\,\mathsf{B}_{\{f,g\}}(w),\\
  \mathsf{B}_f(z)\mathsf{\Theta}_\rho(w)
  &\sim
  K^{(2)}_\Delta(z,w)\,\mathsf{\Theta}_{f\cdot\rho}(w),\\
  \mathsf{\Theta}_\rho(z)\mathsf{\Theta}_\sigma(w)
  &\sim 0,
\end{aligned}
\]
where \((f\cdot\rho)(g)=-\rho(\{f,g\})\).  The multi-index singular
products are
\[
  \bigl(\mathsf{J}_x\bigr)_{(0,0)}\mathsf{J}_y
  =
  \mathsf{J}_{[x,y]_{\mathfrak g}},
  \qquad
  \bigl(\mathsf{J}_x\bigr)_{(a,b)}\mathsf{J}_y=0
  \quad\text{for }(a,b)\ne(0,0),
\]
with descendant products obtained by applying the holomorphic
translation operators:
\[
  \partial_z^\alpha\mathsf{J}_x(z)\,\mathsf{J}_y(w)
  \sim
  (-1)^{|\alpha|}\alpha!\,
  \rho^\Delta_{\alpha_1,\alpha_2}(z,w)\,
  \mathsf{J}_{[x,y]_{\mathfrak g}}(w).
\]
\end{thm}

\begin{proof}
The CE factorization algebra is generated by the compactly supported
Dolbeault local Lie algebra
\[
  \Omega_c^{0,\bullet}(B)\widehat\otimes\mathfrak g .
\]
For disjoint polydisks the product is extension by zero.  When two
insertions collide, the only singular operation in the CE differential
is the local Lie bracket
\([x,y]_{\mathfrak g}\); no second singular binary operation exists in
the shifted-cotangent Lie algebra.  The Dolbeault resolution of the
diagonal identifies the singular quotient of the product of two
holomorphic insertions with
\[
  H^2_\Delta(\mathcal O_{\C^2\times\C^2})\,d\zeta_1d\zeta_2,
\]
and the normalized generator is \(K^{(2)}_\Delta\), since its residue on
\(1\) is \(1\).  Therefore the collision of the two currents is
\(K^{(2)}_\Delta\mathsf{J}_{[x,y]}\).

Substituting
\[
  [f,g]=\{f,g\},\qquad
  [f,\rho[1]]=(f\cdot\rho)[1],\qquad
  [\rho[1],\sigma[1]]=0
\]
gives the three displayed OPEs.  The coefficient extraction formula
follows from
\[
  \operatorname{Res}_\Delta(\rho^\Delta_{a,b}\zeta_1^m\zeta_2^n)
  =
  \delta_{a,m}\delta_{b,n}.
\]
The descendant identity is
\[
  \partial_{z_1}^{\alpha_1}\partial_{z_2}^{\alpha_2}
  \frac{1}{\zeta_1\zeta_2}
  =
  (-1)^{|\alpha|}\alpha_1!\alpha_2!\,
  \zeta_1^{-\alpha_1-1}\zeta_2^{-\alpha_2-1}.
\]
Associativity of all iterated two-dimensional OPEs is the Jacobi
identity in \(\mathfrak g\), and the compatibility with \(Q\) is the
Chevalley--Eilenberg identity \(d_{\mathrm{CE}}^2=0\).
\end{proof}

\begin{cor}[No hidden higher holomorphic OPE operations]
\label{cor:no-hidden-higher-two-dimensional-ope}
The complete \(n\)-point holomorphic OPE of the formal Hamiltonian BF
factorization algebra is generated by
Theorem~\ref{thm:formal-two-dimensional-holomorphic-ope}, the regular
commutative product, and the Chevalley--Eilenberg differential.  There is
no primitive \(\ell_n\)-type singular operation for \(n\ge3\).
\end{cor}

\begin{proof}
The coefficient algebra is the CE algebra of the strict dg Lie algebra
\(\Omega_c^{0,\bullet}(B)\widehat\otimes\mathfrak g\).  Its only
nonlinear structure map is the binary Lie bracket of \(\mathfrak g\);
the arity-\(n\) collision maps are therefore the iterated binary
collision maps with the Koszul signs of the CE differential.  Jacobi is
the equality of the two binary collision orders along the small
diagonal.  This proves both associativity and the absence of independent
higher singular coefficients.
\end{proof}

\begin{cor}[Quantum and brane trace OPE]
\label{cor:quantum-brane-trace-two-dimensional-ope}
After Moyal quantization of the Hamiltonian pair and scalar reduction on
the brane,
\[
  \mathsf{B}_f(z)\mathsf{B}_g(w)
  \sim
  K^{(2)}_\Delta(z,w)
  \left(
    \mathsf{B}_{\{f,g\}_\hbar}(w)
    +
    \hbar N\,\omega(f,g)\mathbf 1
  \right),
\]
where
\[
  \{f,g\}_\hbar=\hbar^{-1}(f*g-g*f)\bmod \C\cdot1,\qquad
  \omega(f,g)=[\{f,g\}]_0 .
\]
For \(f=z_1\) and \(g=z_2\),
\[
  \mathsf{B}_{z_1}(z)\mathsf{B}_{z_2}(w)
  \sim
  \hbar N\,K^{(2)}_\Delta(z,w)\mathbf 1 .
\]
This is the Capelli scalar as a two-dimensional holomorphic contact OPE.
\end{cor}

\begin{proof}
The degree-zero quantum brane algebra replaces the Poisson bracket by
the Moyal commutator.  The scalar-reduced Hamiltonian quotient removes
the constant term on the closed side, while the Chan--Paton trace sends
\(1\) to \(N\).  Hence the constant Taylor coefficient of the classical
bracket contributes the central cocycle \(\omega(f,g)\), and the
quantum normalization multiplies it by \(\hbar N\).  For
\((f,g)=(z_1,z_2)\), \(\omega(z_1,z_2)=1\), giving the displayed
contact term.
\end{proof}

\begin{rmk}[Physical reading]
The free BF contraction is
\[
  \alpha_f(z)\,\beta_\rho(w)
  \sim
  \langle\rho,f\rangle K^{(2)}_\Delta(z,w),
\]
with \(\alpha\) the Hamiltonian field and \(\beta\) the shifted
cotangent field.  The cubic interaction
\[
  I_{\mathrm{BF}}
  =
  \frac12\int_{\C^2}\langle\beta,\{\alpha,\alpha\}\rangle\,dz_1dz_2
\]
turns the free contraction into the current OPE of
Theorem~\ref{thm:formal-two-dimensional-holomorphic-ope}.  The quantum
brane trace adds precisely the centre-of-mass scalar of
Corollary~\ref{cor:quantum-brane-trace-two-dimensional-ope}.
\end{rmk}
